\section{PromptNode-Augmented Multi-Agent Graphs}
\label{sec:method}

We model a multi-agent reasoning workflow as a directed communication graph whose nodes may be ordinary agents, tools, memory stores, verifiers, or typed virtual PromptNodes. A PromptNode is not a post-hoc answer patch. It is a first-class graph element with an identity, read/write permissions, a forwarding contract, and lifecycle actions. In our implementation, PromptNode identities include Memory, Instruction-Verifier, Tool-Resource, Verifier, and Agent-like roles. PromptEdges carry explicit read, write, and forward constraints: they determine what evidence a downstream node may consume, what state may be written, and whether an answer label may be forwarded to the final aggregator.

The graph policy follows a two-stage organization. First, diagnostic PromptNodes write structured evidence such as the operative instruction, option binding, definition anchors, adversarial checks, contradictions, confidence, and local memory summaries. Second, a conservative gate decides whether the candidate answer should be forwarded or pruned. The gate is trained or fixed only over graph-visible features, never over hidden labels, task identifiers, or post-hoc correctness. Lifecycle events are recorded as update, compress, prune, merge, and upgrade actions; R419 replays 121 lifecycle events from the fresh R413 DefinitionChallenge run.

For STEM tasks, PromptNodes instantiate InstructionVerifierPN, ToolResourcePN, OptionVerifierPN, and SafetyGatePN. For non-mathematical concept-definition tasks, the DefinitionChallenge path uses InstructionAmbiguityVerifierPN, ConceptDefinitionAnchorPN, OptionSemanticConsistencyVerifierPN, AdversarialVerifierPN, ContradictionDetectorPN, and ConsensusConflictResolverPN. This separation is deliberate: the method is graph organization, while the node library is task-family specific.
